
\section{Using \elan}

The language can be used either through an interpreter or a compiler
as described below in section~\ref{Run the interpreter}
and \ref{how_to_compile} respectively. 


%%%%%%%%%%%%%%%%%%%%%%%%%%%%%%%%%%%%%%%%%%%%%%%%%%%%%%%%%%%%
\subsection{How to run the \elan{} interpreter}
\label{Run the interpreter}
%%%%%%%%%%%%%%%%%%%%%%%%%%%%%%%%%%%%%%%%%%%%%%%%%%%%%%%%%%%%
You can now execute \elan{} by just typing the \verb|elan| command with the
following usage: 
{\footnotesize\begin{verbatim}

        ******** ELAN version 3.6
        (c) LORIA (CNRS, INPL, INRIA, UHP, U-Nancy 2)
            1994, 1995, 1996, 1997, 1998, 1999, 2000, 2001

Usage:
        elan [options] lgi_file [spc_file]
Options:
        --dump, -d              : dump of signature, strategies and rules
        --trace, -t [num]       : tracing of execution (default max)
        --statistic, -s/-S      : statistics: short/long
        --warningsoff, -w       : suppress warning messages of the parser
        --quiet, -q             : quiet regime of execution
        --batch, -b             : batch regime (no messages at all)
        --elanlib               : elanlib
        --secondlib, -l lib     : second elan library (by default ..)
        --command, -C           : command language
        --export, --cexport     : export to a .ref file
        --import                : import from a .ref file

\end{verbatim} }

\noindent
Both the logic description file (\verb|lgi_file|) and  
the specification file (\verb|spc_file|)
may be written without their suffixes. 
Default suffixes for \elan{} files are:
\begin{description}
\item[\tt .lgi] -- for the top level logic description file,
\item[\tt .eln] -- for modules used in the top level logic description file,
\item[\tt .spc] -- for the specification file,
\item[\tt .ref] -- for the file in \REF{} format.
\end{description}
%
The system \elan{} searches these files 
(i.e. \verb|.eln|, \verb|.lgi|, \verb|.spc|, \verb|.ref|)
in the following directories with the decreasing priorities:
\begin{itemize}
  \item in the current directory,
  \item in the directory described by the value of a 
    system variable {\tt SECONDELANLIB},
  \item in the directories described as {\tt ELANLIB/commons},
    {\tt ELANLIB/noquote}, {\tt ELANLIB/strategy} and {\tt ELANLIB/ref},
    in this order.
\end{itemize}
%
The default value of the variable {\tt ELANLIB} is set to the subdirectory {\tt elan\_library/elanlib} of the source directory but it can be locally replaced by the switch \verb|--elanlib|. 
The default value of the variable {\tt SECONDELANLIB} is set to the
parent directory, however, it can be also locally replaced by the
switch \verb|--secondlib| (or, \verb|-l|) (for details, see below).

In Chapter \ref{short}, we have shown, how to simply run the 
\elan{} interpreter with a logic description. We now illustrate
on several examples the usage of the different switches.

\begin{description}
%-------------------------------
 \item[{\tt -d}, (or, {\tt --dump})]
      dumps out the signature of the described logic
      (i.e. all sorts and function symbols with their
      profiles), definitions of all strategies and rewrite rules.
      For the example \verb|poly3| from Chapter \ref{short}:
{\footnotesize\begin{verbatim}
> elan -d poly3 someVariables
\end{verbatim} }
\noindent
we obtain a list of rewrite rules of the form:
{\footnotesize\begin{verbatim}      
local rule expand:poly/poly3[Vars] [1/3/fsym=233/vars=9]
 ( VAR(0)+(( VAR(1)+ VAR(2))*( VAR(3)+ VAR(4))))         => 
                     ( VAR(0)+( VAR(5)+( VAR(6)+( VAR(7)+ VAR(8)))))
     where VAR(8) := (simplify:poly/poly3[Vars]) ( VAR(2)* VAR(4))
     where VAR(7) := (simplify:poly/poly3[Vars]) ( VAR(2)* VAR(3))
     where VAR(6) := (simplify:poly/poly3[Vars]) ( VAR(1)* VAR(4))
     where VAR(5) := (simplify:poly/poly3[Vars]) ( VAR(1)* VAR(3))

local rule expand:poly/poly3[Vars] [2/4/fsym=233/vars=6]
 ( VAR(0)+(( VAR(1)+ VAR(2))* VAR(3)))   => ( VAR(0)+( VAR(4)+ VAR(5)))
     where VAR(5) := (simplify:poly/poly3[Vars]) ( VAR(2)* VAR(3))
     where VAR(4) := (simplify:poly/poly3[Vars]) ( VAR(1)* VAR(3))
. . . . . . . . . . . . . . . . . . . . . . . . . . . . . . etc
\end{verbatim} }
\noindent 
Names of rewrite rules (e.g. \verb|expand|) are displayed
together with their sorts (e.g. \verb|poly|), over which these rules
work, and with a module name, where they have been defined
(e.g. \verb|poly3[Vars]|).  
An internal code of the top-most function symbol of its left-hand side 
and the number of variables of each rule are also displayed as an
additional information.  
Variables appearing in rewrite rules are displayed in a uniform way: \verb|VAR(i)|, where
$i$ is an internal index.

\noindent
The print-out continues with a list of strategies in the form:
{\footnotesize\begin{verbatim}
strategy global simplify:poly/poly3[Vars]
 repeat* (
   dc one(expand:poly )
 );
 repeat* (
   dc one(factorize:poly )
 )
. . . . . . . . . . . . . . . . . . . . . . . . . . . . . . etc
\end{verbatim} }
%
%{\footnotesize\begin{verbatim}
%strategy global last_simplify:poly/poly3[Vars]
% repeat* (
%   dc(expand:poly )
% );
% repeat* (
%   dc(factorise:poly )
% )
%
%strategy global simplify:poly/poly3[Vars]
% repeat* (
%   dc(expand:poly );
%   dc(factorise:poly );
%   dc(factorise:poly );
%   repeat* (
%     dc(factorise:poly )
%   )
% )
%. . . . . . . . . . . . . . . . . . . . . . . . . . . . . . etc
%\end{verbatim} }
\noindent The label of the strategy is composed of the name of the
strategy, its sort and the name of the module where it has been
defined (as for labelled rules).

The specification of the signature of the described logic 
is composed of a set of defined (and internal) sorts:
{\footnotesize\begin{verbatim}
Sorts
intern int, ident, identifier, string, list[identifier], 
intern ident, int, variable, intern string, bool, 
poly, 
\end{verbatim} }
and of a list of function symbols with their profiles in the following form:
{\footnotesize\begin{verbatim}
Function symbols
 @                     : (variable)poly        pri 0 code 231;
 @                     : (int)poly             pri 0 code 232;
 @ '+' @               : (poly poly)poly       assocRight      pri 1 code 233;
'(' @ '+' @ ')'        : (poly poly)poly       pri 0 code 233;
 @ '*' @               : (poly poly)poly       assocRight      pri 2 code 234;
'(' @ '*' @ ')'        : (poly poly)poly       pri 0 code 234;
'deriv' '(' @ ',' @ ')': (poly variable)poly   pri 0 code 235;
. . . . . . . . . . . . . . . . . . . . . . . . . . . . . . etc
\end{verbatim} }
%-------------------------------
 \item[{\tt -t} (or, {\tt --trace})] allows tracing \elan{}
programs.  The maximal depth of tracing can be specified, 
like \verb|-t 9|, however, by default (i.e. \verb|-t|), it traces all levels.  The
number specified as the trace level counts matchings of a left-hand
side, verifications of the conditional part of a rule, and evaluations
of strategies applied over terms in local affectations. Thus, if the
user wants to trace his/her rewrite derivation up to the 5-th level, the
trace argument should be specified as \verb|-t 15|.
%
{\footnotesize\begin{verbatim} > elan -t poly3 someVariables
. . . . . . . . . . . . . . . . . . . . . . . . . . . . . . etc enter
query term finished by the key word 'end':
deriv(X*X,X) end

[] start with term :
    deriv((X*X),X)

       [reduce] start:
       [0]  deriv((X*X),X)
                   [reduce] start:
                   [0] (X*deriv(X,X))
                   [1] (X*1)
                   [2]  X
                   [reduce] stop :
               applying strategy 'simplify:poly/poly3[Vars]' on  X
                 trying  'expand:poly' on        X
                 fail of 'expand:poly' 
                 trying  'expand:poly' on        X
                 fail of 'expand:poly' 
                 trying  'factorize:poly' on     X
                 fail of 'factorize:poly' 
                 trying  'factorize:poly' on     X
                 fail of 'factorize:poly' 
                 trying  'factorize:poly' on     X
                 fail of 'factorize:poly' 
             setting VAR(4) on  X
                   [reduce] start:
                   [0] (deriv(X,X)*X)
                   [1] (1*X)
                   [2]  X
                   [reduce] stop :
               applying strategy 'simplify:poly/poly3[Vars]' on  X
                 trying  'expand:poly' on        X
                 fail of 'expand:poly' 
                 trying  'expand:poly' on        X
                 fail of 'expand:poly' 
                 trying  'factorize:poly' on     X
                 fail of 'factorize:poly' 
                 trying  'factorize:poly' on     X
                 fail of 'factorize:poly' 
                 trying  'factorize:poly' on     X
                 fail of 'factorize:poly' 
             setting VAR(5) on  X
                   [reduce] start:
                   [0] (X+X)
                   [reduce] stop :
               applying strategy 'simplify:poly/poly3[Vars]' on (X+X)
                 trying  'expand:poly' on       (X+X)
                 fail of 'expand:poly' 
                 trying  'expand:poly' on       (X+X)
                 fail of 'expand:poly' 
                 trying  'factorize:poly' on    (X+X)
                 fail of 'factorize:poly' 
                 trying  'factorize:poly' on    (X+X)
                 fail of 'factorize:poly' 
                 trying  'factorize:poly' on    (X+X)
                 fail of 'factorize:poly' 
             setting VAR(3) on (X+X)
       [1] (X+X)
       [reduce] stop :
     trying  'expand:poly' on   (X+X)
     fail of 'expand:poly' 
     trying  'expand:poly' on   (X+X)
     fail of 'expand:poly' 
     trying  'factorize:poly' on        (X+X)
     fail of 'factorize:poly' 
     trying  'factorize:poly' on        (X+X)
     fail of 'factorize:poly' 
     trying  'factorize:poly' on        (X+X)
     fail of 'factorize:poly' 

[] result term:
   (X+X)

[] end
\end{verbatim} }

%-------------------------------
 \item[{\tt -s} ({\tt -S}, or {\tt --statistics})]
   displays
   a brief ({\tt -s}) or a complete ({\tt -S}) version of statistics. These
   statistics contain the information about the running time
   of the last query, average speed in rewrites per second,
   the statistics of labelled and unlabelled rules, which were tried and applied.
{\footnotesize\begin{verbatim}
> elan -S poly3 someVariables
 . . . . . . . . . . . . . . . . . . . . . . . . . . . . . . etc
 enter query term finished by the key word 'end':
 deriv(X*X,X) end

[] start with term :
    deriv((X*X),X)

[] result term:
   (X+X)

[] end

 Statistics:
 total time     (0.007+0.000)=0.007 sec (main+subprocesses)

 average speed  714 inf/sec
                5 nonamed rules applied,    23 tried
                0   named rules applied,    20 tried
 named rules
           applied   tried    rule for symbol
                 0       8    expand:poly
                 0      12    factorize:poly
 nonamed rules
           applied   tried    rule for symbol
                 0       4    ( poly + poly )
                 2      12    ( poly * poly )
                 3       7    deriv( poly , variable )
 end of statistics
\end{verbatim} } 
%
The complete statistics show a detailed list of
applications and tries for labelled and unlabelled rules.
The brief statistics do not precise the last paragraph about 
unlabelled rules.

%-------------------------------
 \item[{\tt -w} (or, {\tt --warningsoff})]
   eliminates all ambiguity warnings
   produced during parsing of the \elan{} description of a logic.

 \item[{\tt -q} (or, {\tt --quiet})]
   suppresses all error messages and warnings during the
   execution of the \elan{} program.

 \item[{\tt -b} (or, {\tt --batch})]
   suppresses all messages, thus this mode
   is useful when the \elan{} system is 
   executed from a script file in the batch mode.

%-------------------------------
 \item[{\tt --elanlib dir}]
   locally redefines the value
   of the system variable {\tt ELANLIB} by the string \verb|dir|.

 \item[{\tt --secondlib dir} (or, {\tt --secondlib dir})]
   does the same with
   the variable {\tt SECONDELANLIB}.

%-------------------------------
 \item[{\tt --export fname}, {\tt --cexport fname}]
  exports the input specification into the \REF{} format stored in the file {\tt fname}.
  The switch {\tt cexport} is used only when the produced \REF{} file is compiled by the
        \elan{} compiler.

 \item[{\tt --import fname}] 
  imports a specification in the \REF{} format from the \REF{} file {\tt fname}
  produced, in general, by the switch {\tt --export}.
{\footnotesize\begin{verbatim}
> elan --export poly3.ref poly3 someVariables.spc

        ******** ELAN version 3.6 ********

         (c)   LORIA (CNRS, INPL, INRIA, UHP, U-Nancy 2), 1994 - 2001
 . . . . . . . . . . . . . . . . . . . . . . . . . . . . . . etc

Export file poly3.ref has been created

> elan --import poly3.ref
        ******** ELAN version 3.6 ********

         (c)   LORIA (CNRS, INPL, INRIA, UHP, U-Nancy 2), 1994 - 2001

Import form file poly3

Importing Identifiers
Importing Sorts
Importing Modules
Importing RuleNames
Importing StrategyNames
Import rule extractrule1:identifier/list[identifier]!LO form module list[identifier]
Import rule extractrule2:identifier/list[identifier]!LO form module list[identifier]
Import rule expand:poly/poly3[Vars]!LO form module poly3[Vars]
Import rule expand:poly/poly3[Vars]!LO form module poly3[Vars]
Import rule factorize:poly/poly3[Vars]!LO form module poly3[Vars]
Import rule factorize:poly/poly3[Vars]!LO form module poly3[Vars]
Import rule factorize:poly/poly3[Vars]!LO form module poly3[Vars]
Import strategy listExtract:identifier/list[identifier]
Import strategy simplify:poly/poly3[Vars]
Import strategy last_simplify:poly/poly3[Vars]

enter query term finished by the key word 'end':
\end{verbatim} }

%-------------------------------
 \item[{\tt -C}]
   runs the \elan\ interpreter
   with a simple command language interpreted on the top-most level.
   This mode is described in Section~\ref{TopLevelInter}.

%-------------------------------
% \item[
%   {\tt -c} ({\tt --compiler}),
%   {\tt -O} ({\tt --optimise}),
%   {\tt --code},
%   {\tt -o} ({\tt --output}),
%   {\tt -n} ({\tt --deterministic}) and
%   {\tt --exe}] concern the \elan{} compiler and they will
%   be explained in Section~\ref{how_to_compile}.
\end{description}

\noindent
After loading all modules specified in the logic description file, 
\elan{} requires to enter a query term, which
has to be finished by the word 'end' or by character \verb|^D| (control-D).
Its evaluation can be interrupted by typing \verb|^C|, 
the interpreter then proposes a simple run-time menu:
\begin{itemize}
  \item {\tt ExecutionAbort} - A -
    aborts the current execution and allows to enter another query term,
  \item {\tt Continue} - C -
    continues the interrupted execution,
  \item {\tt Dump} - D -
    dumps the signature, rules and strategies of the currently loaded
    logic,
  \item {\tt Exit} - E -
    quits the \elan{} system,
  \item {\tt Statistics} - S, s -
    writes the actual statistics corresponding to the last evaluated query,
  \item {\tt ChangeTrace} - T -
    changes trace level of the current derivation such that 
    the interpreter asks for a new trace level,
  \item {\tt ChangeQuiet} - Q -
    switches between two displaying modes: quiet -- unquiet,
\comment{
  \item {\tt Input} - I -
    changes the input sort of the query,
  \item {\tt Output} - O -
    changes the output sort of the query,
  \item {\tt stRategy} - R -
    changes the name of the top-most strategy applied on the entered query.
}
\end{itemize}


%%%%%%%%%%%%%%%%%%%%%%%%%%%%%%%%%%%%%%%%%%%%%%%%%%%%%%%%%%%%
\subsection{Top level interpreter} \label{TopLevelInter}
%%%%%%%%%%%%%%%%%%%%%%%%%%%%%%%%%%%%%%%%%%%%%%%%%%%%%%%%%%%%

The command language of the \elan\ interpreter improves the user's interface 
by offering to the user several commands, which allow to:
\begin{itemize}
  \item change several parameters of the system,
  \item debug programs using  break-points,
  \item keep the history of queries and their results,
  \item runs the current logic with different strategies, different entries, etc.
\end{itemize}

\noindent
When the \elan{} interpreter is started with the option \verb|-C|, instead of asking:\\
\verb|enter query term finished by the key word 'end'| \\
\elan{} prompts the user: \\
\verb|enter command finished by ';'| \\
We show a brief description of all commands of the top level command language
%(also found in the file {\tt ELANLIB/help.txt}), 
which will be later illustrated on several examples.
%
\begin{description}
  \item[{\tt quit}]   
    terminates the session and leaves the  \elan{} interpreter.
%  \item[{\tt help}]   
%    shows the help file {\tt ELANLIB/help.txt}.
  \item
    [{\tt load mod}]  
    loads (i.e. imports) an additional module, where the identifier \verb|mod| 
    describes its name (possibly with its arguments),
  \item
    [{\tt batch mod}] 
    calls a script file with the name 
    \verb|mod| (an identifier), which contains a stream of commands
    of the command language.
    It tries to evaluate all commands of the script file, and
    it returns the control to the interactive mode.
    The script files can be also concatenated.
 \item
   [{\tt qs}]  shows the queue of input queries (i.e. the history up to the
       last ten queries).
       The terms in this queue of queries {\tt qs} can be
       referenced by the identifiers
       \verb|Q|, \verb|Q1|, ..., \verb|Q9| respecting their sorts.
       These sorted identifiers could be used in constructions of
       further queries.
 \item
   [{\tt rs}] shows the queue of the latest results 
         (i.e. the history up to the last ten results).
       Terms in this queue of results {\tt rs} are
       referenced by identifiers
       \verb|R|, \verb|R1|, ..., \verb|R9|.
       In the case of non-deterministic computations, 
       there is no correspondence between query terms
       \verb|Qi| and result terms \verb|Ri|.
  \item
    [{\tt startwith (str)term}]  
    redefines the 'start with' pattern originally defined in the 
    logic description file \verb|.lgi|.
    The pattern \verb|term| may contain also symbols
          \verb|query| standing as a place-holder of an input term
          entered by the user using the command \verb|run|,
          or symbols \verb|Q|, $\dots$, \verb|Qi| or 
          \verb|R|, $\dots$, \verb|Ri|
          standing for values of previously put queries or results.
  \item
    [{\tt checkwith term}]
    redefines the 'check with' pattern originally defined in the 
    logic description file.
    The boolean pattern \verb|term| may contain any symbol amongst
    \verb|query|, 
    \verb|Q|, $\dots$, \verb|Qi| or 
          \verb|R|, $\dots$, \verb|Ri|.
  \item
    [{\tt printwith term}]
    sets-up the `print-with' pattern. The `print-with' pattern
    may contain the symbol \verb|result|, which is a place-holder for
    result terms. Intuitively, the meaning of the `print-with' term $P$
    is that any result $r$ of the computation is substituted for the
    place-holder \verb|result| in the pattern $P$. The substituted term
    $P[result/r]$ is normalized and printed out.
    This command can be used for the automatic transformation
    of results into different formats, e.g.
    a transformation of 
    very large result terms into more readable or incomplete
    notation, or a conversion into a \LaTeX\ notation.
    The default value of the `print-with' pattern is \verb|result|, i.e.
        the identity.
  \item
    [{\tt run term}]
    evaluates the term \verb|term| as a query w.r.t. rewrite rules of the current logic 
    and eventually returns results.
    The entry term \verb|term| may contain identifiers
    \verb|Q|, \verb|Q1|, ..., \verb|Q9|,
    (resp. \verb|R|, \verb|R1|, ..., \verb|R9|) 
    referring to values of preceding queries or results, e.g.
        items of the queue of queries {\tt qs}
    (resp. results {\tt rs}).
  \item
    [{\tt sorts type type type}]  
    sets-up sorts of the symbols \verb|query|, 
    \verb|result| and of the `print-with' pattern.
  \item
    [{\tt stat}]
    prints the statistics,
  \item
    [{\tt dump}] 
    prints rules, strategies and the signature,
  \item
    [{\tt dump name}]
    dumps only labelled rules with the name \verb|name|,
    strategies labelled by \verb|name| or
    unlabelled rules with the head symbol \verb|name|.
  \item
    [{\tt dump n}]     
    dumps only unlabelled rules of a function symbol with the 
    internal code \verb|n|.
  \item
    [{\tt display n}]  
    if \verb|n| is not zero, all printed terms are in the internal form.
    The internal form of a term attaches to a symbol name its internal code,
    which allows to identify overloaded function symbols.
    {\tt display 0} switches all print-outs into traditional (thus, overloaded)
    format.
  \item 
    [{\tt trace n}]
     sets-up the trace level to the value \verb|n|.
  \item
    [{\tt break name}] (resp. {\tt unbreak name})
    sets (resp. removes) break-points to all labelled rules 
    with the name \verb|name|, to all strategies 
    labelled by \verb|name| and to all unlabelled rewrite rules with 
    the head symbol \verb|name|.
  \item
    [{\tt break n}] (resp. {\tt unbreak n})
    sets (resp. removes) break-points to unlabelled rules with the head symbol 
    with the internal code \verb|n|.
  \item 
    [{\tt breaks}]
    shows a list of all breakpoints.
\end{description}

\noindent
We illustrate the command language on a small \elan{} session with
the example \verb|poly3| from the previous Chapter.
When the switch \verb|-C| is specified, the \elan{} interpreter  loads
the module \verb|Query| describing the syntax of commands.
{\footnotesize\begin{verbatim}
> elan -q -C poly3 someVariables

        ******** ELAN version 3.6 ********

         (c)   LORIA (CNRS, INPL, INRIA, UHP, U-Nancy 2), 1994 - 2001
   handling module identifier
     handling module ident
       handling module bool
       end of bool
     end of ident
     handling module eq[identifier]
       handling module cmp[identifier]
       end of cmp[identifier]
     end of eq[identifier]
   end of identifier
   handling module list[identifier]
     handling module int
       handling module builtinInt
       end of builtinInt
     end of int
   end of list[identifier]
   handling module poly3[Vars]
     handling module eq[variable]
       handling module cmp[variable]
       end of cmp[variable]
     end of eq[variable]
   end of poly3[Vars]
 handling module Query[poly,poly,poly]
 end of Query[poly,poly,poly]
\end{verbatim} }

\noindent
We can run simple queries using the command \verb|run|:
{\footnotesize\begin{verbatim}
enter command finished by ';':
run deriv(X*X,X);   
[] start with term :
    deriv((X*X),X)
[] result term:
   (X+X)
[] end

enter command finished by ';':
run deriv(X*X*X+X*X+X+1,X);
[] start with term :
    deriv(((X*(X*X))+((X*X)+(X+1))),X)
[] result term:
   (1+(X*(((X*2)+2)+X)))
[] end
\end{verbatim} }

\noindent
The queries and their results are stored in two queues, which can be
listed by commands \verb|qs|, \verb|rs|:
{\footnotesize\begin{verbatim}
enter command finished by ';':
qs;
Queries ... 2 : elements 
Q1      poly     deriv((X*X),X)
Q       poly     deriv(((X*(X*X))+((X*X)+(X+1))),X)

enter command finished by ';':
rs;
Results ... 2 : elements 
R1      poly    (X+X)
R       poly    (1+(X*(((X*2)+2)+X)))
\end{verbatim} }

\noindent
The previous queries and their results can be 
referenced in the construction of the new ones:
{\footnotesize\begin{verbatim}
enter command finished by ';':
run deriv(R,X) ;
[] start with term :
    deriv((1+(X*(((X*2)+2)+X))),X)
[] result term:
   ((X*6)+2)
[] end

enter command finished by ';':
run deriv(Q*Q,X) ;
[] start with term :
    deriv((deriv((1+(X*(((X*2)+2)+X))),X)*deriv((1+(X*(((X*2)+2)+X))),X)),X)
[] result term:
   (24+(X*72))
[] end
\end{verbatim} }

\noindent
Function symbols defined by unlabelled rules can be debugged, 
for example, 
if we know their internal code. In the following example, the command
\verb|dump| displays the symbol with the internal code \verb|216|,
and then, we can put a break-point on this function.
{\footnotesize\begin{verbatim}
dump 235;
function dump
'deriv' '('  @ ','  @ ')'        : (poly variable)poly   pri 0 code 235;

enter command finished by ';':
break 235;
function dump
'deriv' '('  @ ','  @ ')'        : (poly variable)poly   pri 0 code 235;
\end{verbatim} }

\noindent
Having break-points set, 
the execution shows all entry and exit points of traced functions, or
strategies:
{\footnotesize\begin{verbatim}
enter command finished by ';':
run deriv(X*X,X);

[] start with term :
    deriv((X*X),X)
 [0]  deriv((X*X),X)
       [0]  deriv(X,X)
       [1]  1
       [0]  deriv(X,X)
       [1]  1
 [1] (X+X)
[] result term:
   (X+X)
[] end
\end{verbatim} }

\noindent
We can switch all outputs into the internal format by the command 
\verb|display 1|:
{\footnotesize\begin{verbatim}
enter command finished by ';':
display 1;

enter command finished by ';':
run deriv(X*X,X);
[] start with term :
    deriv_235((X_228*X_228),X_228)
 [0]  deriv_235((X_228*X_228),X_228)
       [0]  deriv_235(X_228,X_228)
       [1]  1
       [0]  deriv_235(X_228,X_228)
       [1]  1
 [1] (X_228+X_228)
[] result term:
   (X_228+X_228)
[] end
\end{verbatim} }

\noindent
We can change the `start-with' pattern to \verb|(last_simplify)deriv(query,X)|,
and the command \verb|stat| informs us about the current setting and the
performance of the last executed query.
{\footnotesize\begin{verbatim}
enter command finished by ';':
startwith (last_simplify)deriv(query,X);

enter command finished by ';':
stat;
Input type : poly
Output type: poly
Print type : poly
Strategy   : last_simplify
Start_with :  deriv_235( VAR(0),X_228)
Check_with :  true_1
Print_with :  VAR(0)
------------

 Statistics:
 total time     (0.013+0.010)=0.023 sec (main+subprocesses)

 average speed  384 inf/sec
                5 nonamed rules applied,    23 tried
                0   named rules applied,    20 tried
 named rules
           applied   tried    rule for symbol
                 0       8    expand:poly
                 0      12    factorize:poly
 nonamed rules
           applied   tried    rule for symbol
                 0       4    ( poly + poly )
                 2      12    ( poly * poly )
                 3       7    deriv( poly , variable )
 end of statistics
\end{verbatim} }




%%%%%%%%%%%%%%%%%%%%%%%%%%%%%%%%%%%%%%%%%%%%%%%%%%%%%%%%%%%%
\subsection{How to run the \elan{} compiler}
\label{how_to_compile}
%%%%%%%%%%%%%%%%%%%%%%%%%%%%%%%%%%%%%%%%%%%%%%%%%%%%%%%%%%%%
The \elan{} compiler transforms a logic description into an executable
binary file.  The compiler (for detailed description,
see~\cite{Vittek-RTA96,MoreauK-PLILP+ALP98}) produces an efficient
\verb|C| code, which is later compiled by the \verb|cc| or the
\verb|gnu gcc| compiler. 
The executable binary code is dependent on the particular
architecture, because a small part of the \elan{} library performing
the basic non-deterministic operations is written in assembly
language~\cite{Moreau-JICSLPW98}.
Up to now, this is the only reason that makes
the compiler available only on the architectures DEC-ALPHA, SUN4 and
Intel-PC.  This release of the \elan{} compiler does not compile the
whole \elan{} language: the main restriction is that the \elan{}
compiler does not treat input/output and reflective capabilities of
the \elan{} language. There is also a restriction on the class of rules
containing AC-symbols that can be compiled.
Thus, any logic description containing AC-symbols should be carefully
studied before being compiled (see Section~\ref{compilo} for more
details).\\


\noindent To run the \elan\ compiler, a Java Virtual Machine (JAVA 2 SDK, aka
jdk 1.2, or newer), a C compiler (gcc recommended) and a C-shell (tcsh
recommended) have to be installed on your system. 
To complete the installation, the script {\tt PWD/bin/elanc} has to be
modified: {\tt CLASSPATH} and {\tt JAVA} variables have to be set according to
your system. 
Then, you can now execute the \elan\ compiler by just typing the \verb|elanc|
command with the following usage: 
{\footnotesize\begin{verbatim}
elanc [[elan options][compiler options]] file[.lgi] [file.[spc]]
Options:
        -output <name>
        -verbose
        -nocode
        -nosplit
        -quiet
        -debug
        -noOptimiseChoicePoint [default]
        -optimiseChoicePoint
        -proofterm
\end{verbatim}
}

\noindent
By default, the compiler produces a lot of \verb|C|~files and a
\texttt{lgiName.makefile} file, which are later compiled.
\begin{description}
\item [The switch] \verb|-output nqueens| produces the executable
  file with the name \verb|nqueens|. 

\item [The switch] \verb|-nosplit| generates only two files:
  \texttt{lgiName.c} and \texttt{lgiName.core.c}. This switch is
  useful to compile small specifications.

\item [The switch] \verb|-nocode| directly builds an executable file and
  deletes all generated \verb|C|~files. 

\item [The switch] \verb|-quiet| suppresses all printed messages during
  compilation phase.

\item [The switch] \verb|-debug| activates the debugging mode: more
  messages are printed during compilation and the generated code becomes
  less efficient.

\item [The switch] \verb|-noOptimiseChoicePoint [default]| avoids the 
  deterministic analysis phase.

\item [The switch] \verb|-optimiseChoicePoint| reduces the number of
  set choice-points: this switch reduces the runtime needed memory and
  improve the efficiency.
  It is recommended to not activate this switch if your specification
  contains AC symbols.

\item [The switch] \verb|-proofterm| generates the complete trace of the
  normalisation process with respect to a set of labelled rules by the
  strategy \rw{normin} described in Section~\ref{language}.
%\item -color [default]
%\item -noColor
%\item -choicePointDebug
%\item -oldVariableAffectation
%\item -strategy <0|1|2>
%\item -withGoto
\end{description}

\noindent
We illustrate the compiler on a simple example of a distributed file.
The first example shows the compilation of the program \verb|nqueens|
with the query \verb|queens(8)| (given in the \texttt{.lgi} file):

{\footnotesize\begin{verbatim}
> elanc -output queens -nosplit -nocode -optimiseChoicePoint nqueens
elan --cexport nqueens.ref ./nqueens.lgi

        ******** ELAN version 3.6 ********

         (c)   LORIA (CNRS, INPL, INRIA, UHP, U-Nancy 2), 1994 - 2001
   handling module nqueens
     handling module int
       handling module builtinInt
         handling module bool
         end of bool
       end of builtinInt
     end of int
   end of nqueens

Export file nqueens.ref has been created

java -classpath classes.zip:elanlib/Compiler/classes REM nqueens.ref
-output queens -nosplit -nocode -optimiseChoicePoint
Reduce ELAN Code Machine. Reading from file nqueens.ref . . .
Parsing...........
Compiling nonamed rules...................
Compiling strategies..
Building queens...

> queens
result = .([](4),.([](2),.([](7),.([](3),.([](6),.([](8),.([](5),.([](1),nil))))))))
result = .([](5),.([](2),.([](4),.([](7),.([](3),.([](8),.([](6),.([](1),nil))))))))
result =
.([](3),.([](5),.([](2),.([](8),.([](6),.([](4),.([](7),.([](1),nil))))))))
...
after 91 solutions:
result = .([](5),.([](7),.([](2),.([](6),.([](3),.([](1),.([](4),.([](8),nil))))))))

rewrite_step = 227356
\end{verbatim} }

\noindent
It is also possible to get the whole trace of the execution. In order
to produce this information the executable code has to be compiled in
the debugging mode and the switch {\tt -trace} has to be used:

{\footnotesize\begin{verbatim}
> elanc -debug -nosplit nqueens
 ...
Reduce ELAN Code Machine. Reading from file nqueens2.ref . . .
Parsing
Compiling nonamed rules...................
Compiling strategies..
execute the following line to build your program
        gmake -f nqueens.make
> gmake -f nqueens.make
gcc -pipe -g3 -DDEBUG -DPDEBUG -DBITSET32 ...
...

> a.out -trace 
|start with: fun_221_queens()_([](8))
|rewrite[1] queens(,,)([](8),[](8),nil)
|start with: str_176(queens(,,)([](8),[](8),nil))
| start with: fun_211_greater_int(,)_([](8),[](0))
| rewrite[2] true
| start with: str_278([](8))
| rewrite(str_278)[3] [](1)
| start with: fun_222_noattack(,,)_([](1),[](1),nil)
| rewrite[4] true
| start with: fun_202_minus(,)_([](8),[](1))
| rewrite[5] [](7)
| start with: str_176(queens(,,)([](7),[](8),.([](1),nil)))
|  start with: fun_211_greater_int(,)_([](7),[](0))
|  rewrite[6] true
|  start with: str_278([](8))
|  rewrite(str_278)[7] [](1)
|  start with: fun_222_noattack(,,)_([](1),[](1),.([](1),nil))
|   start with: fun_210_neq_int(,)_([](1),[](1))
|   rewrite[8] false
|  rewrite[8] noattack(,,)([](1),[](1),.([](1),nil))
...
\end{verbatim}}

\noindent
A more detailed description of the compiler can be found in
Section~\ref{compilo}. 

\Attention{
Another way to call the \elan\ compiler (also known as {\em
  Reduced \elan\ Machine}) consists in running Java as follows:\\
\texttt{java -classpath \$\{ELANLIB\}/Compiler/classes REM} \nt{REF file name}
\nt{options}\\
The \nt{REF file name} must be created by the option \texttt{--cexport}
\nt{REF file name} of the \elan\ interpreter simply used here to parse the 
\nt{LGI file name}. 
}


%%%%%%%%%%%%%%%%%%%%%%%%%%%%%%%%%%%%%%%%%%%%%%%%%%%%%%%%%%%%
\section{Keep in touch}
%%%%%%%%%%%%%%%%%%%%%%%%%%%%%%%%%%%%%%%%%%%%%%%%%%%%%%%%%%%%

In case you have any problems, questions or comments about \elan{}, 
or just only to be keeping in touch with the \elan{} team,
please use the following e-mail address:\\
\indent\mbox{\tt elan@loria.fr}
or post on the elan-users mailing-list:\\
\indent\mbox{\tt elan-users@loria.fr}

\noindent
More informations on \elan\ can be found on the web page:\\
\indent\localisationWebElan.

\noindent
\elan\ is issued from the PROTHEO research team, information about
this research project can be found at the following address:\\
\indent\verb|http://www.loria.fr/equipes/protheo|.

